% =============================================================================
% Section 4: Analysis and Evaluation
% For: Constraint Acquisition of Configuration Knowledge Bases with MSS
% =============================================================================

\section{Analysis and Evaluation}
\label{sec:evaluation}

\textsc{ConGen} and \textsc{AcqMss} have the following properties.

% -----------------------------------------------------------------------------
\subsection{Time Complexity}
% -----------------------------------------------------------------------------

The \textit{worst-case} complexity of \textsc{AcqMss} in terms of the number of needed consistency
checks for determining one maximum satisfiable subset $B'$ is $2\gamma \times \log_2(\frac{n}{\gamma}) + 2\gamma$ where
$\gamma$ is the \textit{number of elements deleted from} $B$, $n$ is the \textit{number of constraints} in $B$, and
$2\gamma$ represents the branching factor and the number of leaf-node consistency checks. In
this context, we interpret the checking of all positive examples of $E^+$ for consistency
with $NE \cup B \cup BG$ as \textit{one consistency check}. The corresponding \textit{best-case} complexity
is $\log_2(\frac{n}{\gamma}) + 2\gamma$. In the worst case, each element to be excluded from $B$ is located in a
different path of the search tree. The factor $\log_2(\frac{n}{\gamma})$ represents the depth of a path of
the \textsc{AcqMss} search tree. In the best case, each to be excluded element is located in
the same path of the \textsc{AcqMss} search tree.

If $n$ is the number of constraints in $B$ removed by \textsc{AcqMss}, \textsc{QuickXPlain}~\cite{junker2004}
requires $\log_2(\frac{n}{\pi}) + 2\pi$ consistency checks in the best case and $2\pi \times \log_2(\frac{n}{\pi}) + 2\pi$ in the
worst case ($\pi$ is the assumed conflict set size). With an increasing number of conflicts
in $B$, traditional HSDAG-based conflict resolution~\cite{reiter1987} performance deteriorates.

% --- Theorems ---
\begin{theorem}[Correctness of \textsc{AcqMss}]
\label{thm:correctness}
Let $B' \subseteq B$ (the constraint bias) be the set of
constraints returned by \textsc{AcqMss} -- then $B'$ accepts all positive examples.
\end{theorem}

\begin{proof}
\textsc{AcqMss} is only activated, if the following property is fulfilled: $\forall e^{+}_i \in E^{+}$:
$consistent(\{e^{+}_i\} \cup NE \cup BG)$. If fulfilled, \textsc{AcqMss} incrementally aggregates the sets $B'_{\beta}$ into $BG$
which fulfill: $\forall e^+ \in E^+$: $consistent(\{e^+\} \cup NE \cup BG)$. This way, all $e^+ \in E^+$ are accepted. \qed
\end{proof}

\begin{theorem}[Completeness of \textsc{ConGen}]
\label{thm:completeness}
If $\forall e^+_i \in E^+$ : $consistent(\{e^+_i\} \cup NE \cup BG)$,
\textsc{ConGen} will return $B' \subseteq B$.
\end{theorem}

\begin{proof}
If $\forall e^+_i \in E^+$: $consistent(\{e^+_i\} \cup NE \cup BG)$, \textsc{ConGen} activates \textsc{AcqMss} which
returns $B' \subseteq B$. In the worst case, all constraints are removed, i.e., $B' = \emptyset$. \qed
\end{proof}

\begin{remark}
If $B' = \emptyset$, \textsc{ConGen} will only return the negated negative examples $NE$ derived from $E^-$.
\end{remark}

\begin{corollary}
\textsc{AcqMss} fails to compute a bias reduction $B'$ if and only if the following
condition is violated: $\exists e^+_i \in E^+$: $inconsistent(\{e^+_i\} \cup NE \cup BG)$.
\end{corollary}

% -----------------------------------------------------------------------------
\subsection{Experimental Setup}
\label{sec:experimental_setup}
% -----------------------------------------------------------------------------

All experiments were conducted on a MacBook Pro with Apple M4 Pro, 48\,GB RAM, running macOS~15. All algorithms were implemented on top of the FLAMA 2.0.1 framework~\cite{Galindo2023flama}. The SAT solver used is Glucose4 via the PySAT library~\cite{imms-sat18} in incremental mode. Random seeds are fixed per experiment to ensure reproducibility; all seeds and configurations are available in the public artifact\footnote{Artifact URL omitted for double-blind review.}.

\paragraph{Bias generation.}
For each feature model, the constraint bias $B$ is generated by enumerating all admissible binary constraints over the feature vocabulary $(V,D)$ using the constraint language $L = \{x_i \circ x_j \mid i \neq j,\; \circ \in \{\wedge, \neg\wedge, \rightarrow\}\}$. Each constraint is translated into one or more CNF clauses for SAT-based reasoning. Constraints that encode the \textit{feature hierarchy} (mandatory, optional, or-group, alternative-group relationships) are placed in the background knowledge~$BG$ and are not subject to learning. The remaining constraints -- specifying cross-tree requirements (\textit{requires}) and incompatibilities (\textit{excludes}) -- form the elements of the constraint bias~$B$.

\paragraph{Cross-validation protocol.}
We employ \textit{3-fold cross validation}: the example set $E = E^+ \cup E^-$ is split into three disjoint folds of approximately equal size. In each iteration, two folds serve as the training set and the remaining fold as the test set. The constraint bias ordering is shuffled per fold to mitigate potential ordering effects in \textsc{AcqMss}. We report the mean and standard deviation of accuracy across the three folds.

\paragraph{Comparison strategies.}
To evaluate the quality of a learned knowledge base~$KB$ against the ground-truth constraint theory~$C_\tau$ (derived from the original feature model), we employ three complementary comparison strategies:
\begin{enumerate}
    \item \textbf{Description-based} (primary): compares the human-readable constraint descriptions (e.g., ``\textit{id requires db}'') in the learned $KB$ against those in $C_\tau$. A constraint is a true positive if its description appears in both $KB$ and $C_\tau$. This is the most intuitive strategy and is used for the accuracy results reported in the main tables.
    \item \textbf{Clause-based} (structural): converts both $KB$ and $C_\tau$ to normalized CNF clause sets and computes set-level precision, recall, and F1-score. This captures structural equivalence at the clause level.
    \item \textbf{Semantic} (SAT-based): checks bidirectional entailment between $KB$ and $C_\tau$ via SAT solving. A constraint in $C_\tau$ is entailed if $KB \cup BG \models c$, and vice versa. This is the strongest notion of equivalence and accounts for logically equivalent but syntactically different representations.
\end{enumerate}

% -----------------------------------------------------------------------------
\subsection{Evaluation Dataset}
% -----------------------------------------------------------------------------

In our evaluation, we apply real-world configuration knowledge bases~\cite{heradio2022} for
the purpose of (oracle-based) query answering. These knowledge bases are represented in
terms of feature hierarchies and corresponding constraints between individual features.
Positive (and negative) examples can stem from (non-)solutions generated by the knowledge base.
In such a scenario, a solution generator creates different configurations (examples) which are then evaluated by the oracle
(i.e., does the generated configuration satisfy the knowledge base?).
We kept constraints specifying the feature
hierarchy in the background knowledge ($BG$) and constraints specifying component
requirements and incompatibilities as elements of the constraint bias $B$.

Table~\ref{tab:fm_summary} summarizes the three feature models used in our evaluation, spanning different scales from a small IDE plugin configurator (14~features) to a medium-sized software product line (179~features).

\begin{table}[htbp]
\centering
\caption{Summary of feature models used in the evaluation.}
\label{tab:fm_summary}
\begin{tabular}{lrrrl}
\toprule
 & \#features & $|B|$ & \#clauses & domain \\
\midrule
$KB_1$ (REAL-FM-7) & 14 & 295 & 314 & IDE plugins \\
$KB_2$ (FQA)        & 179 & 459 & 932 & quality attributes \\
$KB_3$ (Arcade)     & 65 & 1{,}755 & 1{,}960 & arcade game SPL \\
\bottomrule
\end{tabular}
\end{table}

% -----------------------------------------------------------------------------
\subsection{Sampling Methods}
% -----------------------------------------------------------------------------

For example generation, we use the feature model as an oracle and apply three feature selection
strategies frequently used in testing scenarios~\cite{pereira2021}.
\textit{Random sampling} (RS) generates examples with the total number equal to $n$, $2n$, $3n$, and $m$ where $n$ is the
number of features and $m$ is the smallest number of valid configurations
including each pair of features.
\textit{$n$-wise coverage} ($n$-COV) produces a set of configurations satisfying coverage heuristics; \textit{2-wise coverage} (2-COV) ensures
each combination of two features is included at least once.
\textit{Feature frequency} (FF) generates examples such that each
feature is at least once included and excluded. For each generated configuration, the oracle
classifies it as positive or negative.

Table~\ref{tab:example_sizes} shows the resulting example set sizes.

\begin{table}[htbp]
\centering
\caption{Number of positive ($|E^+|$) and negative ($|E^-|$) examples per sampling strategy.}
\label{tab:example_sizes}
\begin{tabular}{lrrrrrr}
\toprule
 & \multicolumn{2}{c}{$KB_1$} & \multicolumn{2}{c}{$KB_2$} & \multicolumn{2}{c}{$KB_3$} \\
\cmidrule(lr){2-3} \cmidrule(lr){4-5} \cmidrule(lr){6-7}
Strategy & $|E^+|$ & $|E^-|$ & $|E^+|$ & $|E^-|$ & $|E^+|$ & $|E^-|$ \\
\midrule
RS($n$)  &  13 &  1 & 162 & 17 &  59 &  6 \\
RS($2n$) &  26 &  2 & 323 & 35 & 117 & 13 \\
RS($3n$) &  38 &  4 & 484 & 53 & 176 & 19 \\
RS($m$)  &   8 &  1 &  15 &  1 &  13 &  1 \\
2-COV    &   0 &  9 &   0 & 16 &   1 & 13 \\
FF       &   6 &  3 &  59 &  7 &  17 &  5 \\
\bottomrule
\end{tabular}
\end{table}

A notable observation is that the 2-COV strategy generates predominantly (or exclusively) negative examples for all three knowledge bases. This is because 2-wise coverage systematically exercises all pairwise feature combinations, many of which violate the feature model constraints.

% -----------------------------------------------------------------------------
\subsection{Evaluation of Generated Knowledge Bases}
% -----------------------------------------------------------------------------

The example set is evaluated using \textit{3-fold cross validation}. Each generated knowledge base
(constraint theory) is tested in terms of \textit{accuracy} (see Formula~\ref{eq:accuracy}).

\begin{equation}
\label{eq:accuracy}
accuracy(KB) = \frac{TP + TN}{TP + TN + FP + FN}
\end{equation}

Accuracy is measured on the basis of \textit{true positives} ($TP$) (number of positive examples accepted by $KB$), \textit{true negatives} ($TN$) (number of rejected negative examples),
\textit{false positives} ($FP$) (number of accepted negative examples), and \textit{false negatives} ($FN$)
(number of rejected positive examples).

% --- Table: Performance (#checks and runtime) ---

\subsubsection{Performance}

Table~\ref{tab:performance} depicts the average number of \textsc{AcqMss} consistency checks and the average runtime (in milliseconds) needed for determining one maximal satisfiable subset (MSS).

\begin{table}[htbp]
\centering
\caption{\textsc{AcqMss} \#consistency checks~/~runtime (in \textit{msec}) for different real-world configuration knowledge bases.}
\label{tab:performance}
\begin{tabular}{lrrr}
\toprule
Strategy & $KB_1$ & $KB_2$ & $KB_3$ \\
\midrule
RS($n$)  & 515~/~717   & 706~/~12{,}228   & 3{,}090~/~81{,}297 \\
RS($2n$) & 538~/~1{,}500  & 713~/~34{,}426   & 3{,}241~/~181{,}130 \\
RS($3n$) & 544~/~2{,}304  & 744~/~67{,}512   & 3{,}290~/~301{,}179 \\
RS($m$)  & 466~/~425   & 476~/~841     & 2{,}187~/~14{,}738 \\
2-COV    &  10~/~113   &  10~/~1{,}043    & 1{,}429~/~3{,}594 \\
FF       & 452~/~355   & 622~/~3{,}657    & 2{,}743~/~22{,}964 \\
\bottomrule
\end{tabular}
\end{table}

The number of consistency checks and runtime scale with both the bias size $|B|$ and the number of examples. $KB_3$ ($|B|{=}1{,}755$) requires up to $6{\times}$ more consistency checks than $KB_1$ ($|B|{=}295$). The 2-COV strategy achieves the fewest consistency checks across all knowledge bases (only 10 for $KB_1$ and $KB_2$), because it generates predominantly negative examples, resulting in few constraints being removed from the bias. As the number of random samples increases from RS($n$) to RS($3n$), both consistency checks and runtime increase due to larger training sets.

% --- Accuracy tables ---

\subsubsection{Accuracy of \textsc{ConGen}}

Table~\ref{tab:accuracy_rs} shows the average accuracy for random sampling strategies.

\begin{table}[htbp]
\centering
\caption{\textit{Accuracy} of generated knowledge bases with \textit{random sampling} (RS).}
\label{tab:accuracy_rs}
\begin{tabular}{lrrr}
\toprule
Strategy & $KB_1$ & $KB_2$ & $KB_3$ \\
\midrule
RS($n$)  & 0.278 $\pm$ 0.048 & 0.854 $\pm$ 0.069 & 0.415 $\pm$ 0.037 \\
RS($2n$) & 0.558 $\pm$ 0.218 & 0.927 $\pm$ 0.043 & 0.509 $\pm$ 0.134 \\
RS($3n$) & 0.860 $\pm$ 0.143 & 0.967 $\pm$ 0.006 & 0.554 $\pm$ 0.066 \\
RS($m$)  & 0.194 $\pm$ 0.174 & 0.367 $\pm$ 0.153 & 0.194 $\pm$ 0.174 \\
\bottomrule
\end{tabular}
\end{table}

Table~\ref{tab:accuracy_2cov} shows the accuracy with 2-wise coverage.

\begin{table}[htbp]
\centering
\caption{\textit{Accuracy} of generated knowledge bases with \textit{2-wise coverage} (2-COV).}
\label{tab:accuracy_2cov}
\begin{tabular}{lr}
\toprule
 & Accuracy \\
\midrule
$KB_1$ & 0.556 $\pm$ 0.385 \\
$KB_2$ & 1.000 $\pm$ 0.000 \\
$KB_3$ & 0.944 $\pm$ 0.096 \\
\bottomrule
\end{tabular}
\end{table}

Table~\ref{tab:accuracy_ff} shows the accuracy with feature frequency.

\begin{table}[htbp]
\centering
\caption{\textit{Accuracy} of generated knowledge bases with \textit{feature frequency} (FF).}
\label{tab:accuracy_ff}
\begin{tabular}{lr}
\toprule
 & Accuracy \\
\midrule
$KB_1$ & 0.444 $\pm$ 0.193 \\
$KB_2$ & 0.786 $\pm$ 0.078 \\
$KB_3$ & 0.222 $\pm$ 0.048 \\
\bottomrule
\end{tabular}
\end{table}

Table~\ref{tab:accuracy_all} consolidates the accuracy results across all strategies.

\begin{table}[htbp]
\centering
\caption{Consolidated accuracy for all sampling strategies and knowledge bases.}
\label{tab:accuracy_all}
\begin{tabular}{lrrr}
\toprule
Strategy & $KB_1$ & $KB_2$ & $KB_3$ \\
\midrule
RS($n$)  & 0.278 & 0.854 & 0.415 \\
RS($2n$) & 0.558 & 0.927 & 0.509 \\
RS($3n$) & 0.860 & 0.967 & 0.554 \\
RS($m$)  & 0.194 & 0.367 & 0.194 \\
2-COV    & 0.556 & 1.000 & 0.944 \\
FF       & 0.444 & 0.786 & 0.222 \\
\bottomrule
\end{tabular}
\end{table}

The accuracy results reveal several observations.
First, \textbf{more examples lead to higher accuracy}: for random sampling, accuracy consistently improves from RS($m$) to RS($3n$) across all knowledge bases. For $KB_2$ (FQA), accuracy increases from $0.367$ with RS($m$) to $0.967$ with RS($3n$), demonstrating the importance of sufficient training data.

Second, \textbf{2-wise coverage (2-COV) achieves the highest accuracy for medium and large knowledge bases}: $KB_2$ achieves perfect accuracy ($1.000$) and $KB_3$ achieves $0.944$. The 2-COV strategy systematically covers all pairwise feature interactions, which is particularly effective for knowledge bases with complex inter-feature constraints.

Third, \textbf{the smallest knowledge base ($KB_1$) is more sensitive to sampling strategy}: with only 14 features, $KB_1$ shows high variance across strategies ($0.194$--$0.860$). The small number of features limits the diversity of generated examples, making the choice of sampling strategy more critical.

Fourth, \textbf{RS($m$) consistently underperforms} because the minimum number of examples is often insufficient to capture the constraint structure, especially for larger knowledge bases.

% --- Bias reduction ---

\subsubsection{Bias Reduction}

Table~\ref{tab:kb_size} provides insight into the bias reduction achieved by \textsc{AcqMss} and \textsc{Reduce}.

\begin{table}[htbp]
\centering
\caption{Bias reduction: average $|MSS|$ (maximal satisfiable subset size) and $|KB|$ (final knowledge base size after \textsc{Reduce}) per strategy.}
\label{tab:kb_size}
\begin{tabular}{lrrrrrr}
\toprule
 & \multicolumn{2}{c}{$KB_1$ ($|B|$=295)} & \multicolumn{2}{c}{$KB_2$ ($|B|$=459)} & \multicolumn{2}{c}{$KB_3$ ($|B|$=1{,}755)} \\
\cmidrule(lr){2-3} \cmidrule(lr){4-5} \cmidrule(lr){6-7}
Strategy & $|MSS|$ & $|KB|$ & $|MSS|$ & $|KB|$ & $|MSS|$ & $|KB|$ \\
\midrule
RS($n$)  &  92 & 16.7 & 225 & 131.7 &  541 & 173.7 \\
RS($2n$) &  69 & 17.0 & 216 & 134.3 &  371 & 240.7 \\
RS($3n$) &  62 & 12.3 & 201 & 136.0 &  314 & 221.3 \\
RS($m$)  & 132 & 18.3 & 344 & 108.7 & 1{,}167 &  75.7 \\
2-COV    & 294 &  4.0 & 458 & 106.0 & 1{,}375 &  46.0 \\
FF       & 142 & 19.0 & 278 & 124.0 &  817 &  95.3 \\
\bottomrule
\end{tabular}
\end{table}

\textsc{AcqMss} achieves substantial constraint elimination. For $KB_1$, the final $|KB|$ ranges from $4.0$ (2-COV) to $19.0$ (FF), representing a reduction of $93.6\%$--$98.6\%$ from the original bias. An interesting pattern emerges: strategies with predominantly negative examples (2-COV) retain larger MSS values (almost the entire bias remains satisfiable with positive examples) but produce smaller final KBs after the \textsc{Reduce} step, because the negated negative examples ($NE$) subsume many bias constraints. Conversely, strategies with more positive examples (RS($3n$)) produce smaller MSS but may retain more constraints in the final KB.

% --- Comparison with iterative approaches ---

\subsubsection{Comparison with Iterative Approaches}

We compare \textsc{ConGen} (passive learning) with an iterative (active) constraint acquisition approach inspired by QuAcq~\cite{bessiere2013}. In the iterative approach, examples are processed one at a time, and after each example the already learned constraint theory is taken into account when selecting the next query. We evaluate two query selection modes:

\begin{itemize}
    \item \textbf{Example-only}: queries are drawn exclusively from the pre-generated example pool. For each candidate example, the algorithm checks two conditions: (1)~the example must be consistent with the already learned KB and background knowledge, and (2)~the example must violate at least one constraint still in the remaining bias. This mode stops when the pool is exhausted or no suitable example remains.
    \item \textbf{Example-first}: queries are first drawn from the example pool (as in example-only mode). When the pool is exhausted or no pool example satisfies the two conditions, the algorithm falls back to \textit{SAT-based query generation}: for each remaining bias constraint~$c_i$, a SAT solver searches for a configuration that satisfies $KB \cup BG$ but violates~$c_i$. This mode continues until no more queries can be generated or a maximum query limit is reached.
\end{itemize}

The key difference is that \textit{example-first} can generate additional queries beyond the provided examples by leveraging SAT solving, whereas \textit{example-only} is limited to the given example pool. In both modes, the oracle (feature model) provides ground-truth classification.

Table~\ref{tab:iterative_accuracy} compares the test-fold accuracy, and Table~\ref{tab:iterative_semantic} compares the \textit{semantic F1-score} -- the structural similarity between the learned KB and the ground-truth constraint theory~$C_\tau$.

\begin{table}[htbp]
\centering
\caption{Test-fold \textit{accuracy} comparison: \textsc{ConGen} (passive) vs.\ iterative approaches.}
\label{tab:iterative_accuracy}
\begin{tabular}{llrrr}
\toprule
Mode & Strategy & $KB_1$ & $KB_2$ & $KB_3$ \\
\midrule
\multicolumn{5}{l}{\textit{\textsc{ConGen} (passive)}} \\
 & RS($n$)  & 0.278 & 0.854 & 0.415 \\
 & RS($3n$) & 0.860 & 0.967 & 0.554 \\
 & 2-COV    & 0.556 & 1.000 & 0.944 \\
 & FF       & 0.444 & 0.786 & 0.222 \\
\midrule
\multicolumn{5}{l}{\textit{Iterative -- example-only}} \\
 & RS($n$)  & 0.944 & 0.944 & 0.953 \\
 & RS($3n$) & 0.952 & 0.978 & 0.964 \\
 & 2-COV    & 0.667 & 0.500 & 0.556 \\
 & FF       & 0.889 & 0.940 & 0.958 \\
\midrule
\multicolumn{5}{l}{\textit{Iterative -- example-first}} \\
 & RS($n$)  & 1.000 & 0.961 & 0.953 \\
 & RS($3n$) & 0.974 & 0.978 & 0.964 \\
 & 2-COV    & 0.889 & 0.667 & 0.556 \\
 & FF       & 0.889 & 0.985 & 0.958 \\
\bottomrule
\end{tabular}
\end{table}

\begin{table}[htbp]
\centering
\caption{\textit{Semantic F1-score} (structural KB quality) comparison: \textsc{ConGen} vs.\ iterative approaches.}
\label{tab:iterative_semantic}
\begin{tabular}{llrrr}
\toprule
Mode & Strategy & $KB_1$ & $KB_2$ & $KB_3$ \\
\midrule
\multicolumn{5}{l}{\textit{\textsc{ConGen} (passive)}} \\
 & RS($n$)  & \textbf{0.176} & \textbf{0.945} & \textbf{0.525} \\
 & RS($3n$) & \textbf{0.135} & \textbf{0.947} & \textbf{0.568} \\
 & 2-COV    & 0.078 & \textbf{0.891} & \textbf{0.324} \\
 & FF       & \textbf{0.192} & \textbf{0.911} & \textbf{0.513} \\
\midrule
\multicolumn{5}{l}{\textit{Iterative -- example-only}} \\
 & RS($n$)  & 0.042 & 0.046 & 0.046 \\
 & RS($3n$) & 0.046 & 0.048 & 0.038 \\
 & 2-COV    & 0.057 & 0.034 & 0.034 \\
 & FF       & 0.049 & 0.034 & 0.042 \\
\midrule
\multicolumn{5}{l}{\textit{Iterative -- example-first}} \\
 & RS($n$)  & 0.064 & 0.057 & 0.046 \\
 & RS($3n$) & 0.074 & 0.048 & 0.038 \\
 & 2-COV    & \textbf{0.105} & 0.046 & 0.034 \\
 & FF       & 0.049 & 0.056 & 0.042 \\
\bottomrule
\end{tabular}
\end{table}

The two tables reveal a fundamental trade-off between test-fold accuracy and structural knowledge base quality.

\textbf{Iterative approaches achieve higher test-fold accuracy} ($>0.9$ with random sampling) but produce \textit{minimal} KBs that correctly classify held-out examples without capturing the underlying constraint structure. Their semantic recall is consistently below $0.04$, meaning fewer than $4\%$ of the ground-truth constraints are recovered.

\textbf{\textsc{ConGen} produces structurally superior KBs.} For $KB_2$ (FQA), \textsc{ConGen} achieves a semantic F1 of $0.945$--$0.947$ compared to the iterative approaches' $0.034$--$0.057$ -- a difference of more than one order of magnitude. This is because \textsc{ConGen} performs bias reduction over the \textit{entire} constraint set simultaneously, whereas the iterative approach learns constraints one at a time through individual queries.

\textbf{The 2-COV strategy with \textsc{ConGen} achieves accuracy comparable to iterative methods} for $KB_2$ ($1.000$) and $KB_3$ ($0.944$), while retaining significantly better structural quality (semantic F1 of $0.891$ and $0.324$ respectively). This demonstrates that well-designed sampling strategies can close the accuracy gap while preserving the structural advantages of passive learning.

\textbf{The iterative 2-COV mode underperforms} because predominantly negative examples are filtered out by the pool-selection conditions, leaving insufficient queries for effective learning.

% --- Runtime comparison ---

Table~\ref{tab:runtime_comparison} compares the total runtime.

\begin{table}[htbp]
\centering
\caption{Total runtime (in \textit{msec}) comparison (3-fold CV total).}
\label{tab:runtime_comparison}
\begin{tabular}{llrrr}
\toprule
Mode & Strategy & $KB_1$ & $KB_2$ & $KB_3$ \\
\midrule
\multicolumn{5}{l}{\textit{\textsc{ConGen} (passive)}} \\
 & RS($n$)  & 2{,}157 & 36{,}721 & 243{,}976 \\
 & RS($3n$) & 6{,}925 & 202{,}638 & 903{,}774 \\
 & 2-COV    & 342 & 3{,}132 & 10{,}790 \\
\midrule
\multicolumn{5}{l}{\textit{Iterative -- example-only}} \\
 & RS($n$)  & 1{,}430 & 60{,}907 & 133{,}814 \\
 & RS($3n$) & 2{,}034 & 131{,}883 & 289{,}810 \\
 & 2-COV    & 1{,}563 & 19{,}954 & 3{,}994 \\
\midrule
\multicolumn{5}{l}{\textit{Iterative -- example-first}} \\
 & RS($n$)  & 12{,}723 & 258{,}109 & 1{,}169{,}106 \\
 & RS($3n$) & 9{,}887 & 281{,}698 & 655{,}199 \\
 & 2-COV    & 24{,}401 & 467{,}294 & 595{,}213 \\
\bottomrule
\end{tabular}
\end{table}

\textbf{\textsc{ConGen} with 2-COV is the fastest approach} (342\,ms for $KB_1$, 10{,}790\,ms for $KB_3$) due to minimal consistency checks.
\textbf{The iterative example-first mode is the slowest}, requiring SAT-based query generation for each remaining bias constraint. For $KB_3$ with RS($n$), example-first takes $1{,}169$ seconds vs.\ \textsc{ConGen}'s $244$ seconds -- a $4.8{\times}$ overhead.
\textbf{The iterative example-only mode offers the best accuracy-to-runtime ratio} for random sampling strategies, as it avoids the expensive SAT-based query generation while still benefiting from targeted example selection.

% --- Discussion ---

\subsection{Discussion}

Our evaluation demonstrates that \textsc{ConGen} effectively acquires constraint theories from pre-defined examples using real-world feature model knowledge bases. The key findings are:

\begin{enumerate}
    \item \textbf{Scalability.} \textsc{AcqMss} scales to knowledge bases of different sizes. Even for $KB_3$ with $|B|{=}1{,}755$ constraints, the algorithm completes within reasonable time bounds (${\leq}301$ seconds for the largest test suites).

    \item \textbf{Accuracy vs.\ sampling.} The accuracy of the learned knowledge base is highly dependent on the sampling strategy. Coverage-based strategies (2-COV) outperform random sampling for medium and large knowledge bases, achieving up to perfect accuracy.

    \item \textbf{Passive vs.\ active trade-off.} Iterative (active) approaches achieve higher test-fold accuracy ($>0.9$) but produce \textit{structurally impoverished} KBs: their semantic F1 rarely exceeds $0.06$, recovering fewer than $4\%$ of the ground-truth constraints. \textsc{ConGen} produces KBs that are structurally far closer to the ground truth (semantic F1 up to $0.947$), making it the preferred choice when the goal is to \textit{reconstruct} the underlying constraint theory rather than merely classify examples.

    \item \textbf{Comparison strategy matters.} Description-based comparison underestimates KB quality by up to $2.6{\times}$ compared to semantic comparison. Semantic F1 provides the most faithful assessment of structural KB quality.

    \item \textbf{Example-only vs.\ example-first.} The example-first mode provides marginal accuracy improvements at significant runtime cost ($4$--$10{\times}$ slower). The structural quality gain is also marginal.

    \item \textbf{Bias reduction.} The \textsc{Reduce} step is crucial for obtaining compact KBs, eliminating $>93\%$ of constraints in all evaluated scenarios.
\end{enumerate}

\subsection{Threats to Validity}

\textbf{Internal validity.} We use 3-fold cross validation to mitigate overfitting. The bias ordering is shuffled per fold to reduce ordering effects. Random seeds are fixed per experiment for reproducibility.

\textbf{External validity.} Our evaluation covers three real-world feature models of varying sizes (14--179 features). Larger industrial-scale feature models (e.g., Linux kernel with thousands of features) may exhibit different performance characteristics.

\textbf{Construct validity.} We employ three complementary comparison strategies (description, clause, semantic) to avoid relying on a single metric. The semantic F1-score captures logically equivalent constraints that description-based comparison misses, providing a more faithful assessment of structural KB quality.
